\chapter{Chiffrage du développement}\label{chap:chiffrage}

Nous avons voulu savoir ce qu'aurait coûté ce travail s'il avait été confié à un développeur indépendant
expérimenté, habitué à l'émulation et à Android. Ce chiffrage porte uniquement sur ce que \soft ajoute à ZXTune
entre le 2 octobre et le 24 novembre 2026.

\textbf{Le coût des drivers propriétaires NebulaWeave n'est pas inclus et ne sera pas communiqué dans ce document :
il relève du secret des affaires.}

\section{Pourquoi nous ne comptons pas les lignes}
Une première estimation, fondée sur le nombre de lignes ajoutées au dépôt, donnait plus de 900 jours de travail.
Ce chiffre ne tient pas : une bonne part de ces lignes est du code importé tel quel (ymfm, mt32emu, ares), des
ressources, des tests générés en série, ou du code plus long qu'il n'aurait fallu. Un développeur qui connaît son
sujet écrit moins de code pour le même résultat. Nous avons donc chiffré fonctionnalité par fonctionnalité, en
nous demandant à chaque fois combien de jours il faudrait à une personne compétente pour obtenir le même
comportement, tests compris.

Nous distinguons aussi le code créé du code adapté. Adapter un cœur d'émulation existant, par exemple ajouter à
Ootake la lecture de la table d'ondes de chaque voix, demande surtout de comprendre le cœur ; le code ajouté tient
en quelques centaines de lignes. Intégrer une bibliothèque importée ne demande que de vérifier sa licence, de la
compiler pour Android et de la raccorder.

\section{Détail}
\begin{longtable}{@{}p{0.75\textwidth} r@{}}
\toprule \textbf{Fonctionnalité} & \textbf{Jours} \\ \midrule
Sorties par voix et états des puces dans les moteurs (libvgm, SID, SNES, PC Engine, PCM) & 25 \\
Oscilloscope Vulkan et OpenGL ES, déclenchement & 15 \\
Tableau de bord SID et réponse du filtre & 15 \\
Vues FM et inspecteur & 12 \\
Vue Super Nintendo & 6 \\
Vue des puces PCM (YMW258, SCSP) & 6 \\
Vue PC Engine & 4 \\
Cartes NES et MT-32 & 3 \\
Portage d'ares et intégration SNSF & 6 \\
Choix des cœurs FM, préréglages, étage de sortie & 8 \\
Options SID (voie des samples, mono, filtres) & 4 \\
Listes multiples, favoris, file d'attente, M3U & 8 \\
Refonte de l'interface & 15 \\
Formats audio et correctifs divers & 2 \\
Version bureau (Qt) & 12 \\
Outils d'extraction VGM (Model 2 et 3) & 8 \\
Tests et rendus de contrôle & 10 \\
Manuel & 6 \\ \midrule
Sous-total & 165 \\
Conception, recherche, recette (15~\%) & 25 \\ \midrule
\textbf{Total} & \textbf{190} \\
\bottomrule
\end{longtable}

\section{Coût}
\begin{longtable}{@{}l r r@{}}
\toprule \textbf{Hypothèse} & \textbf{Taux journalier} & \textbf{Coût} \\ \midrule
Basse & 450~€ & 85\,500~€ \\
Médiane & 600~€ & 114\,000~€ \\
Haute & 750~€ & 142\,500~€ \\
\bottomrule
\end{longtable}

Ces montants ne comprennent ni la base ZXTune, développée pendant des années par son auteur, ni les bibliothèques
libres réutilisées sans changement, ni les drivers propriétaires. Ils donnent un ordre de grandeur : à vous de les
ajuster si votre taux journalier ou votre rythme diffèrent des nôtres.
